\subsubsection{Ganancia máxima del Circuito 4}

En el Circuito 4 la salida se mide sobre la resistencia \(R\). Por lo tanto, la función transferencia se obtiene mediante un divisor de tensión en un circuito RLC serie:

\begin{equation}
    H_4(s)=\frac{V_R(s)}{V_{in}(s)}
    =
    \frac{R}{R+sL+\frac{1}{sC}}.
\end{equation}

Normalizando:

\begin{equation}
    \boxed{
    H_4(s)=
    \frac{RCs}{LCs^2+RCs+1}
    }
\end{equation}

En régimen sinusoidal:

\begin{equation}
    H_4(j\omega)=
    \frac{R}
    {R+j\left(\omega L-\frac{1}{\omega C}\right)}.
\end{equation}

Por lo tanto, el módulo de la transferencia es:

\begin{equation}
    |H_4(j\omega)|=
    \frac{R}
    {
    \sqrt{
    R^2+
    \left(\omega L-\frac{1}{\omega C}\right)^2
    }
    }.
\end{equation}

La ganancia máxima se obtiene cuando la parte reactiva total del circuito serie se anula:

\begin{equation}
    \omega L-\frac{1}{\omega C}=0.
\end{equation}

Entonces:

\begin{equation}
    \omega_0=\frac{1}{\sqrt{LC}}.
\end{equation}

En resonancia se cumple:

\begin{equation}
    Z_T(\omega_0)=R.
\end{equation}

Por lo tanto:

\begin{equation}
    H_4(j\omega_0)=\frac{R}{R}=1.
\end{equation}

Entonces, la ganancia máxima ideal es:

\begin{equation}
    \boxed{
    |H_4|_{\max}=1
    }
\end{equation}

En decibeles:

\begin{equation}
    \boxed{
    |H_4|_{\max,\mathrm{dB}}
    =
    20\log_{10}(1)
    =
    0\ \mathrm{dB}
    }
\end{equation}

El Circuito 4 no puede ganar más de \(0\,\mathrm{dB}\), ya que es un circuito pasivo. Para cualquier frecuencia se cumple:

\begin{equation}
    \sqrt{
    R^2+
    \left(\omega L-\frac{1}{\omega C}\right)^2
    }
    \ge R.
\end{equation}

Por lo tanto:

\begin{equation}
    |H_4(j\omega)|\le 1.
\end{equation}

La igualdad se alcanza únicamente en resonancia, cuando \(X_L=X_C\). Por eso, el valor máximo posible de la ganancia es \(0\,\mathrm{dB}\).
